% The settings reference: every key of medorama.toml (medorama/src/config.rs: Config::new for
% the defaults, Config::load for the types and checks) and of vitalstream.toml
% (viewers/vitalstream/rust/src/settings.rs, python/vitalstream/vitalstream_py/settings.py,
% cpp/vitalstream/src/settings.hpp; vitalstream.example.toml), the settings Medorama writes for
% its viewer (medorama/src/viewer.rs), and the options fixed when a program is built
% (docs/building.md; medorama/Cargo.toml, viewers/vitalstream/rust/Cargo.toml,
% viewers/vitalstream/cpp/CMakeLists.txt).
\chapter{Settings reference}
\label{chap:settings-reference}

This chapter lists every setting of Medorama's \file{medorama.toml} and of the VitalStream
viewers' \file{vitalstream.toml}, each with its type, its default and its meaning. It then lists
the options that are fixed when a program is built. It is written for administrators and
integrators. How the settings are used is described in \cref{chap:administration} for Medorama
and in \cref{chap:vitalstream} for the viewers; the command-line options that override some of
them are described in \cref{chap:cli}.

\section{medorama.toml}
\label{sec:settings-reference-medorama}

\paragraph{The file.} The settings live in \file{<home>/medorama.toml}, in the catalog's home
folder (\cref{chap:administration}). \cmd{medorama init} writes the file with a comment above
each key. Medorama reads it when a command or the server starts, so restart the server after a
change.
\begin{itemize}
\item \textbf{Paths} are strings. A \code{\textasciitilde} at the start of a path stands for
  the home folder of the account; a relative path is relative to the catalog's home folder.
\item \textbf{Numbers} are TOML integers, 0 or more. \textbf{Switches} are \code{true} or
  \code{false}. \textbf{Lists} are TOML arrays of strings.
\item \textbf{Every key is checked.} A key that the program does not know is an error
  (\code{medorama.toml: unknown settings: ...}). So are a value of the wrong type
  (\code{medorama.toml: bad value for KEY}), an unknown time zone, a \code{database\_url} that is
  not a PostgreSQL URL or that holds a password, and inconsistent sign-in settings
  (\cref{tab:settings-reference-signin-errors}).
\item \textbf{The environment} overrides three settings: \code{MEDORAMA\_DATABASE\_URL} (a
  non-empty value) overrides \code{database\_url}, \code{MEDORAMA\_DATABASE\_PASSWORD} overrides
  \code{database\_password\_file}, and \code{MEDORAMA\_KEY} overrides \code{key\_file}.
\end{itemize}

\subsection{Folders}
\label{sec:settings-reference-folders}

\apientry{results\_dir}{The results folder}
\begin{lstlisting}[style=signature,language=TOML]
results_dir = "PATH"
\end{lstlisting}
A path. \textbf{Required}; \cmd{init -{}-results} sets it. This is the folder the catalog
writes to: the actions' results go to \file{actions/<action>/r<recording>/j<job>/}, and, unless
\code{index\_dir} names another folder, the recordings' indexes go to \file{index/}. Never choose
a folder that holds recordings.

\apientry{index\_dir}{The index folder}
\begin{lstlisting}[style=signature,language=TOML]
index_dir = "PATH"
\end{lstlisting}
A path. Default: \file{<results\_dir>/index}. This folder holds the recordings' indexes, one
folder per index key (\cref{chap:formats-reference}). A standalone viewer whose \code{cache\_dir}
is this folder opens the recordings that the catalog indexed without indexing them again. The
viewer that the catalog runs always uses it.

\apientry{adapters\_dir}{The adapter folder}
\begin{lstlisting}[style=signature,language=TOML]
adapters_dir = "PATH"
\end{lstlisting}
A path. Default: \file{\textasciitilde/.local/share/medorama/adapters} (Windows:
\code{\%LOCALAPPDATA\%\textbackslash medorama\textbackslash adapters}); \cmd{init} writes the
default out in full. This folder holds the adapters, the programs that read recording formats
that Medorama does not read itself (\cref{chap:addons}). \cmd{medorama addons install} installs
adapters into it.

\apientry{data\_roots}{Where files may be imported from}
\begin{lstlisting}[style=signature,language=TOML]
data_roots = ["PATH", ...]
\end{lstlisting}
A list of paths. Default: \code{[]}, which allows any folder. The files that imported rows name
must lie inside these folders, after symbolic links are resolved; an import refuses a file that
lies elsewhere. Administrators browse only inside these folders in the viewers, and a database
file connected from the page must lie inside one of them.

\apientry{actions\_dirs}{The folders of action plug-ins}
\begin{lstlisting}[style=signature,language=TOML]
actions_dirs = ["PATH", ...]
\end{lstlisting}
A list of paths. Default: \code{["actions"]}, the folder \file{actions/} of the home folder.
Each plug-in action is a subfolder that holds an \file{action.toml} (\cref{chap:actions}).

\apientry{formats\_dirs}{The folders of import formats}
\begin{lstlisting}[style=signature,language=TOML]
formats_dirs = ["PATH", ...]
\end{lstlisting}
A list of paths. Default: \code{["formats"]}, the folder \file{formats/} of the home folder.
Each import format described in TOML is a file \file{<key>.toml}; these formats come in addition
to the built-in ones (\cref{chap:importing}).

\subsection{The database}
\label{sec:settings-reference-database}

\apientry{database}{The SQLite database}
\begin{lstlisting}[style=signature,language=TOML]
database = "PATH"
\end{lstlisting}
A path. Default: \file{medorama.sqlite}, in the home folder. This is the SQLite database file,
used when \code{database\_url} is empty. It holds identifiers, so keep it private.

\apientry{database\_url}{A PostgreSQL database instead}
\begin{lstlisting}[style=signature,language=TOML]
database_url = "postgres://USER@HOST[:PORT]/DATABASE[?sslmode=MODE]"
\end{lstlisting}
A string. Default: empty, which selects the SQLite \code{database}. When it is set, the catalog
lives in a PostgreSQL database (version 14 or later), named by a URL that begins with
\code{postgres://} or \code{postgresql://}. The URL must not hold the password: a URL with one is
refused, whether before the host (\code{user:password@}) or as a parameter
(\code{?password=}). A server on another host must offer TLS unless the URL says
\code{sslmode=disable}, and Medorama verifies the server's certificate and name. Each server
opens at most 32 connections to the database (a request waits for one beyond that), checks a
connection that was idle before using it again, and, unless the URL sets them, gives each
connection 10 seconds to connect, TCP keepalives, and 5 minutes for a statement
(\code{?connect\_timeout=\dots}, \code{keepalives=0}, \code{options=-c statement\_timeout=\dots}
in the URL change them). \code{MEDORAMA\_DATABASE\_URL} overrides this setting.

\apientry{database\_password\_file}{The PostgreSQL password}
\begin{lstlisting}[style=signature,language=TOML]
database_password_file = "PATH"
\end{lstlisting}
A path. Default: none. This file holds the password of \code{database\_url}'s account; its
content is read with surrounding whitespace removed. Make it readable by the catalog's account
only. \code{MEDORAMA\_DATABASE\_PASSWORD} overrides this setting.

\apientry{database\_ca\_file}{Trusted certificate authorities}
\begin{lstlisting}[style=signature,language=TOML]
database_ca_file = "PATH"
\end{lstlisting}
A path. Default: none. This file holds certificate authorities (in PEM form) that Medorama
trusts for the database server's certificate, in addition to the system's: for example a private
authority, or a cloud provider's bundle.

\apientry{key\_file}{The database key}
\begin{lstlisting}[style=signature,language=TOML]
key_file = "PATH"
\end{lstlisting}
A path. Default: none. This file holds the key of an encrypted SQLite database, for builds
with the feature \code{encryption}; its content is read with surrounding whitespace removed. The
key must be set before \cmd{init} (\cref{chap:deployment}). \code{MEDORAMA\_KEY} overrides this
setting. A build without encryption refuses to start when a key is configured, and a PostgreSQL
catalog refuses \code{key\_file}.

\subsection{The server}
\label{sec:settings-reference-server}

\apientry{host}{The address to listen on}
\begin{lstlisting}[style=signature,language=TOML]
host = "ADDRESS"
\end{lstlisting}
A string. Default: \code{"127.0.0.1"}, which accepts connections from this machine only. When
the server listens on any other address than \code{127.0.0.1}, \code{localhost} or \code{::1},
or signs users in through a proxy, without \code{secure\_cookies}, it prints a warning at start. \cmd{serve -{}-host} overrides this
setting. With proxy sign-in, keep this address reachable from the proxy only
(\cref{sec:sign-in-requirements}).

\apientry{port}{The port}
\begin{lstlisting}[style=signature,language=TOML]
port = N
\end{lstlisting}
An integer from 0 to 65535. Default: 8090; \cmd{init -{}-port} writes another. This is the
port of the web page and the API. \cmd{serve -{}-port} overrides it.

\apientry{workers}{Jobs run at the same time}
\begin{lstlisting}[style=signature,language=TOML]
workers = N
\end{lstlisting}
An integer. Default: 2; a value below 1 is read as 1. This is the number of jobs (indexing,
derived series, actions) that the server, and each \cmd{medorama work} process, run at the same
time.

\apientry{auto\_index}{Index new recordings in the background}
\begin{lstlisting}[style=signature,language=TOML]
auto_index = true | false
\end{lstlisting}
A switch. Default: \code{true}. While the server runs, it queues the indexing of new and changed
recordings every 15 seconds.

\apientry{check\_interval\_s}{Re-check the recordings}
\begin{lstlisting}[style=signature,language=TOML]
check_interval_s = SECONDS
\end{lstlisting}
An integer. Default: 900. The server checks every recording this often, to find recordings
that are missing, changed or have lost their index. With 0, it never does; \cmd{medorama check}
still does it on demand.

\apientry{timezone}{The catalog's time zone}
\begin{lstlisting}[style=signature,language=TOML]
timezone = "ZONE"
\end{lstlisting}
An IANA time-zone name. Default: \code{"America/New\_York"}; \cmd{init -{}-timezone} writes
another. Medorama reads the times that imported tables write without an offset in this zone, and
shows times in it; the viewer that the catalog runs uses it too. Times are stored in UTC.

\subsection{The waveform viewer}
\label{sec:settings-reference-viewer}

\apientry{viewer}{The viewer the catalog runs}
\begin{lstlisting}[style=signature,language=TOML]
viewer = "auto" | "cpp" | "rs" | "py" | "none" | "PATH"
\end{lstlisting}
A string. Default: \code{"auto"}, which selects the viewer that \cmd{tools/install-medorama.sh}
recorded in \file{medorama.viewer} next to Medorama's program, or else the first one found of
\cmd{vitalstream-cpp}, \cmd{vitalstream-rs} and \cmd{vitalstream-py}. The values \code{"cpp"},
\code{"rs"} and \code{"py"} select that viewer. With \code{"none"}, the catalog runs no viewer:
\ui{View} and the \ui{VitalStream} tab are grayed out, with the reason. Any other value is the
path of a viewer program.

\apientry{viewer\_url}{A viewer running elsewhere}
\begin{lstlisting}[style=signature,language=TOML]
viewer_url = "URL"
\end{lstlisting}
A string. Default: empty. When it is set, for example to \code{"http://host:8060/"}, the
catalog runs no viewer of its own, and \ui{View} opens that viewer in a new browser tab.

\apientry{qc\_verbose}{Mark repaired data}
\begin{lstlisting}[style=signature,language=TOML]
qc_verbose = true | false
\end{lstlisting}
A switch. Default: \code{false}. When it is \code{true}, the viewer that the catalog runs also
outlines, with a dashed line, the data that a reader repaired, such as data blocks that a file
stored out of order and that the reader put back in order. Breaks are shaded either way.

\subsection{Sign-in and the audit log}
\label{sec:settings-reference-sessions}

These settings choose how users sign in, how long a sign-in lasts, and where a copy of the audit
log goes. \Cref{chap:sign-in} explains the two sign-in methods, the security requirements of
proxy sign-in, and example proxy configurations.

\apientry{sign\_in}{How users sign in}
\begin{lstlisting}[style=signature,language=TOML]
sign_in = "password" | "proxy"
\end{lstlisting}
A string. Default: \code{"password"}. With \code{"password"}, users sign in on Medorama's own
form, with the accounts and passwords that Medorama keeps. With \code{"proxy"}, an
authenticating reverse proxy in front of Medorama (for example oauth2-proxy, Apache httpd with
mod\_auth\_openidc or Shibboleth, or nginx with \code{auth\_request}) signs the user in and names
the user's account in a request header. In both cases, accounts, roles and the PHI permission
live in Medorama (\cref{sec:sign-in-proxy}).

\apientry{proxy\_user\_header}{The header that names the user}
\begin{lstlisting}[style=signature,language=TOML]
proxy_user_header = "HEADER"
\end{lstlisting}
A string. Default: \code{"X-Remote-User"}. Proxy mode only. The request header in which the
proxy sends the user's account name. Medorama compares the whole value with its account names,
ignoring case.

\apientry{proxy\_addresses}{Where the proxy connects from}
\begin{lstlisting}[style=signature,language=TOML]
proxy_addresses = ["ADDRESS", ...]
\end{lstlisting}
A list of IP addresses. Default: \code{["127.0.0.1", "::1"]}. Proxy mode only. Medorama
believes the header of \code{proxy\_user\_header} only on connections from these addresses; on a
connection from anywhere else, it ignores the header and the request is not signed in. Each entry
is one address: host names and ranges are not accepted. When Medorama listens on an IPv6 socket,
a proxy that connects over IPv4 may appear as an IPv4-mapped address, such as
\code{::ffff:10.0.0.5}; list it in that form.

\apientry{proxy\_logout\_url}{Where Sign out leads}
\begin{lstlisting}[style=signature,language=TOML]
proxy_logout_url = "URL"
\end{lstlisting}
A string. Default: empty. Proxy mode only. The address that the page's \ui{Sign out} opens: the
proxy's own sign-out, such as \code{"/oauth2/sign\_out"} for oauth2-proxy or
\code{"/Shibboleth.sso/Logout"} for Shibboleth. When it is empty, the page shows its
signed-out screen.

\paragraph{Errors at start.} Medorama refuses to start when the sign-in settings are
inconsistent, so that a half-made change is never silently ignored. The message follows the
path of \file{medorama.toml}:
\begin{xltabular}{\linewidth}{@{}P{5.0cm}L@{}}
\caption{The sign-in settings' errors at start.}
\label{tab:settings-reference-signin-errors}\\
\toprule
\tabhead Mistake & \tabhead Message \\
\midrule
\endfirsthead
\toprule
\tabhead Mistake & \tabhead Message \\
\midrule
\endhead
\bottomrule
\endfoot
A \code{proxy\_...} setting without \code{sign\_in = "proxy"}, even with an empty value &
  \code{proxy\_user\_header, proxy\_addresses and proxy\_logout\_url need sign\_in = "proxy"} \\
An unknown method & \code{sign\_in: "saml" is not a sign-in method (password, proxy)} \\
A header name with spaces or other invalid characters & \code{proxy\_user\_header: "bad header" is
  not a header name} \\
A host name or a range instead of an IP address & \code{proxy\_addresses: "proxy.example.org" is
  not an IP address} \\
An empty list of addresses & \code{proxy\_addresses: name the address the proxy connects from} \\
\end{xltabular}

\apientry{session\_hours}{How long a sign-in lasts}
\begin{lstlisting}[style=signature,language=TOML]
session_hours = HOURS
\end{lstlisting}
An integer. Default: 12; a value below 1 is read as 1. In password mode, a session ends this long
after its sign-in, and its Privacy Mode choice, kept with the session, ends with it. In proxy
mode, Medorama keeps no session of its own: the proxy or the identity provider sets the
session's lifetime, and this is how long a visit's Privacy Mode choice is kept at most.

\apientry{session\_idle\_minutes}{A sign-in left unused}
\begin{lstlisting}[style=signature,language=TOML]
session_idle_minutes = MINUTES
\end{lstlisting}
An integer. Default: 60. In password mode, a session that has not been used for this long ends;
with 0, only \code{session\_hours} ends it. The page's own checks for new data do not count as
use, so an unattended page signs out. In proxy mode, this is the gap (at least one minute) after
which a user's next request counts as a new visit, with a new \code{login} entry in the audit
log.

\apientry{secure\_cookies}{Cookies over HTTPS only}
\begin{lstlisting}[style=signature,language=TOML]
secure_cookies = true | false
\end{lstlisting}
A switch. Default: \code{false}. When \code{true}, Medorama marks its cookies \code{Secure}, so
that browsers send them over HTTPS only. Set it to \code{true} whenever
users reach Medorama over HTTPS through a reverse proxy, in either sign-in mode
(\cref{chap:deployment}).

\apientry{audit\_export}{A copy of the audit log}
\begin{lstlisting}[style=signature,language=TOML]
audit_export = "PATH"
\end{lstlisting}
A path. Default: none. When it is set, Medorama also appends each seal of the audit log and the
entries it covers to this file, every minute, as JSON lines; the file is readable by its owner
only and synced to disk. Put it on storage that the database's administrators cannot change, so
that it keeps a copy of the log out of their reach (\cref{chap:users-privacy}). An entry that the
database cannot take at once is appended to it at once, as a line with \code{"pending": true}.
\code{GET audit/verify} and \cmd{medorama audit -{}-verify} compare the database's seals with
this file.

\subsection{Settings kept in the database}
\label{sec:settings-reference-db}

Administrators make some settings in the page. Medorama keeps them in the database (the table
\code{settings}), not in \file{medorama.toml}, and writes each change to the audit log.
\begin{xltabular}{\linewidth}{@{}P{3.0cm}P{4.0cm}L@{}}
\caption{The settings made in the Admin tab.}
\label{tab:settings-reference-admin}\\
\toprule
\tabhead Setting & \tabhead Default & \tabhead Where, and meaning \\
\midrule
\endfirsthead
\toprule
\tabhead Setting & \tabhead Default & \tabhead Where, and meaning \\
\midrule
\endhead
\bottomrule
\endfoot
What Privacy Mode withholds & Names, dates of birth, dates (moved back), ages above 89, MRNs,
  account numbers; the text of notes and reports is shown, marked & \ui{Admin} $\rightarrow$
  \ui{Privacy Mode} (\cref{chap:users-privacy}) \\
The catalog's date offset & 0 years, 0 months, 36\,570 days, 0 hours & \ui{Admin} $\rightarrow$
  \ui{Privacy Mode}: each value 0 or more \\
Each patient's own days (\code{spread\_days}) & 365 & \ui{Admin} $\rightarrow$ \ui{Privacy
  Mode}: 30 to 3650 \\
How dates are written & \code{2024-01-31} & \ui{Admin} $\rightarrow$ \ui{Dates and times}: also
  \code{01/31/2024}, \code{31/01/2024} or \code{31.01.2024} \\
The clock & 24-hour & \ui{Admin} $\rightarrow$ \ui{Dates and times}: also 12-hour \\
\end{xltabular}

\section{vitalstream.toml}
\label{sec:settings-reference-vitalstream}

\paragraph{The file.} One settings file serves the three viewers. Each viewer reads the file
named by \opt{-{}-config FILE} (an error if it does not exist), else the first of these files
that exists: \file{viewers/vitalstream/vitalstream.toml}, only for a viewer run from the build
folder of a clone of the repository (git ignores that file); and
\file{\textasciitilde/.config/vitalstream/vitalstream.toml}. A \file{vitalstream.toml} in the
working folder is never read unless \opt{-{}-config} names it: a shared data folder could hold
one that runs programs (\code{adapters\_dir}) or opens the viewer to the network (\code{host},
\code{auth\_file}). On start, the viewer says which settings file it read.
\file{viewers/vitalstream/vitalstream.example.toml} is a commented starting point.
\begin{itemize}
\item \textbf{Sections.} Top-level keys apply to every viewer. The keys of a \code{[py]},
  \code{[rs]} or \code{[cpp]} section apply to that viewer only, and override the top-level ones.
  Each viewer ignores the other viewers' sections.
\item \textbf{Paths} are relative to the folder of the settings file; \code{\textasciitilde}
  stands for the home folder.
\item \textbf{Every key is checked.} An unknown key is an error (\code{unknown settings: ...}),
  and so are a value of the wrong type (\code{bad value for KEY}) and an unknown time zone.
\item \textbf{Command-line options} override the file (\cref{chap:cli}).
\end{itemize}

\apientry{timezone}{The display time zone}
\begin{lstlisting}[style=signature,language=TOML]
timezone = "ZONE"
\end{lstlisting}
An IANA time-zone name. Default: \code{"America/New\_York"}. Option: \opt{-{}-timezone}. Only the
display depends on it; data are stored in UTC.

\apientry{data\_roots}{The folders the viewer may open}
\begin{lstlisting}[style=signature,language=TOML]
data_roots = ["PATH", ...]
\end{lstlisting}
A list of paths. Default: the viewer's root, fixed when it was built
(\cref{sec:settings-reference-build}). Option: \opt{-{}-root DIR}, which may be repeated. These
are the folders that the viewer may list and open, inside its root. A folder that contains the
root is narrowed to the root, and a folder outside the root is left out with a warning: the
settings can narrow the root but never widen it. A path is allowed when the file or folder it
leads to, with every symbolic link resolved, lies inside a data root; a link that leads outside is
refused and not listed. With no data root left, the viewer does not start.

\apientry{max\_records}{Records in a folder}
\begin{lstlisting}[style=signature,language=TOML]
max_records = N
\end{lstlisting}
An integer, 0 or more. Default: 5000. This is the largest number of header files, segments and
layout headers included, below a folder of plain records that the viewer opens as one session; a
folder with more is refused. A WFDB collection is not limited (\cref{chap:recordings}).
\code{max\_records}, \code{max\_depth} and \code{max\_scan\_s} are settings of the viewers only:
Medorama has none, and indexes a recording's folder within their defaults.

\apientry{max\_depth}{Folder levels}
\begin{lstlisting}[style=signature,language=TOML]
max_depth = N
\end{lstlisting}
An integer, 0 or more. Default: 20. This is the largest number of folder levels that the viewer
reads below a folder of plain records.

\apientry{max\_scan\_s}{Time to count the records}
\begin{lstlisting}[style=signature,language=TOML]
max_scan_s = SECONDS
\end{lstlisting}
A number of seconds, 0 or more (an integer or a decimal number). Default: 10. This is the
longest time that the viewer spends counting the records of a folder for its summary, before the
folder opens.

\apientry{require\_password}{Refuse to run without a password}
\begin{lstlisting}[style=signature,language=TOML]
require_password = true | false
\end{lstlisting}
A switch. Default: \code{false}. When it is \code{true}, the viewer refuses to start without a
password (\opt{-{}-no-auth} counts as no password), so a viewer built without password support
never starts. Medorama sets it for the viewer it runs.

\apientry{host}{The address to listen on}
\begin{lstlisting}[style=signature,language=TOML]
host = "ADDRESS"
\end{lstlisting}
A string. Default: \code{"127.0.0.1"}, which accepts connections from this machine only; reach
it from elsewhere through an SSH tunnel. Option: \opt{-{}-host}.

\apientry{allowed\_hosts}{The host names the viewer answers for}
\begin{lstlisting}[style=signature,language=TOML]
allowed_hosts = ["NAME", ...]
\end{lstlisting}
A list of strings. Default: \code{[]}. The host names the viewer answers for besides
\code{localhost}, \code{127.0.0.1}, \code{::1} and \code{host}: the names this machine is
reached by when \code{host} is \code{0.0.0.0}, or a reverse proxy's public name; \code{["*"]}
allows any. A request naming another host is refused (\code{400}), which defeats DNS rebinding
(\cref{sec:viewer-api-checks}).

\apientry{port}{The port}
\begin{lstlisting}[style=signature,language=TOML]
port = N
\end{lstlisting}
An integer from 0 to 65535. Default: 8050 (\cmd{vitalstream-py}), 8060 (\cmd{vitalstream-rs}),
8070 (\cmd{vitalstream-cpp}); the example file sets them in the viewers' sections. Option:
\opt{-{}-port}.

\apientry{index\_workers}{Recordings indexed at the same time}
\begin{lstlisting}[style=signature,language=TOML]
index_workers = N
\end{lstlisting}
An integer, 0 or more. Default: 2. This is the number of recordings that the viewer indexes at
the same time.

\apientry{cache\_dir}{The index folder}
\begin{lstlisting}[style=signature,language=TOML]
cache_dir = "PATH"
\end{lstlisting}
A path. Default: \file{\textasciitilde/.cache/vitalstream}; when the viewer runs from a clone,
\file{viewers/vitalstream/sample\_data/outputs/index}. Option: \opt{-{}-cache-dir}. This is the
index folder. The three viewers share it, since they write the same index, and they share it with
Medorama when it is set to the catalog's \code{index\_dir}.

\apientry{adapters\_dir}{The adapter folder}
\begin{lstlisting}[style=signature,language=TOML]
adapters_dir = "PATH"
\end{lstlisting}
A path. Default: \file{\textasciitilde/.local/share/medorama/adapters} (Windows:
\code{\%LOCALAPPDATA\%\textbackslash medorama\textbackslash adapters}), the folder Medorama uses
too. Option: \opt{-{}-adapters-dir}. This is the
folder of adapters (\cref{chap:addons}).

\apientry{auth\_file}{The password file}
\begin{lstlisting}[style=signature,language=TOML]
auth_file = "PATH"
\end{lstlisting}
A path. Default: \file{\textasciitilde/.config/vitalstream/auth.json}. This is where
\opt{-{}-set-password} keeps the password, as a salted scrypt hash with a random signing key, in a
file readable by its owner only. One password serves the three viewers, \cmd{muse-py serve} and
\cmd{dicom-py serve}. No password is set until one is. A viewer built without password support
ignores an existing password: it never asks for one, and is open to whoever reaches its port.
Signing out adds the sign-in's signature to the file's \code{revoked} list, kept until the
sign-in would have expired.

\apientry{secure\_cookie}{The sign-in cookie over HTTPS only}
\begin{lstlisting}[style=signature,language=TOML]
secure_cookie = true | false
\end{lstlisting}
A switch. Default: \code{false}. When it is \code{true}, the sign-in cookie is marked
\code{Secure}, so that the browser sends it over HTTPS only: set it for a viewer behind an HTTPS
proxy.

\apientry{qc\_verbose}{Mark repaired data}
\begin{lstlisting}[style=signature,language=TOML]
qc_verbose = true | false
\end{lstlisting}
A switch. Default: \code{false}. Option: \opt{-{}-qc-verbose}. With \code{false}, the viewer
does not mark the data that a reader repaired, such as blocks that a file stored out of order and
that the reader put back in order. With \code{true}, it outlines such data with a dashed line and
shows the details on hover. Breaks are shaded either way.

\paragraph{An example.} The following file configures a viewer on a server whose recordings
are in two folders:
\begin{lstlisting}[language=TOML]
timezone = "Europe/Paris"
data_roots = ["/data/icu", "/data/holter"]
cache_dir = "/data/vitalstream-index"
index_workers = 4

[cpp]
port = 8075
\end{lstlisting}

\subsection{The settings Medorama writes for its viewer}
\label{sec:settings-reference-medorama-viewer}

The viewer that the catalog runs reads \file{<home>/viewer/vitalstream.toml}, which Medorama
rewrites every time it starts the viewer. The file sets \code{host = "127.0.0.1"}, a free
\code{port}, \code{cache\_dir} to the catalog's \code{index\_dir}, \code{adapters\_dir} to the
catalog's, \code{auth\_file} to \file{<home>/viewer/auth.json}, the catalog's \code{timezone},
\code{data\_roots} to the catalog's (the whole file system when the catalog has none),
\code{require\_password = true}, and \code{qc\_verbose = true} when the catalog's
\code{qc\_verbose} is. It sets no folder limits, so the viewer keeps their defaults. Do not edit
this file: the next start replaces it.

\section{Options fixed when a program is built}
\label{sec:settings-reference-build}

A few options are not settings: they are fixed in the program when it is built, so that no
setting can widen them. The installers pass them (\cref{chap:installation}); a build by hand sets
them as \cref{tab:settings-reference-build} shows.

\begin{xltabular}{\linewidth}{@{}P{2.9cm}P{5.4cm}L@{}}
\caption{The build options of the programs.}
\label{tab:settings-reference-build}\\
\toprule
\tabhead Program & \tabhead Option & \tabhead Effect \\
\midrule
\endfirsthead
\toprule
\tabhead Program & \tabhead Option & \tabhead Effect \\
\midrule
\endhead
\bottomrule
\endfoot
\cmd{medorama} & Cargo feature \code{encryption}
  (\cmd{cargo build -{}-release -{}-locked -p medorama -{}-features encryption}) & Encrypts the
  database with SQLCipher (OpenSSL is built from source, which needs Perl and \cmd{make}). Off by
  default \\
\cmd{vitalstream-rs} & \code{VITALSTREAM\_ROOT=DIR} in the environment of \cmd{cargo build} &
  The viewer's root: it never opens anything outside \code{DIR} (\code{\textasciitilde} is the
  home folder of whoever runs it). Unset: the whole file system \\
\cmd{vitalstream-rs} & \opt{-{}-no-default-features} (without the Cargo feature
  \code{password}, which is on by default) & No password support: \opt{-{}-set-password} fails, an
  existing password is ignored, and the viewer refuses to run where a password is required \\
\cmd{vitalstream-cpp} & CMake \code{-DVS\_ROOT=DIR} (default empty) & The viewer's root. Empty:
  the whole file system \\
\cmd{vitalstream-cpp} & CMake \code{-DVS\_PASSWORD=OFF} (default \code{ON}) & No password
  support \\
\cmd{vitalstream-py} & \code{VITALSTREAM\_ROOT=DIR} when the package is installed & The viewer's
  root \\
\cmd{vitalstream-py} & \code{VITALSTREAM\_PASSWORD=0} when the package is installed & No
  password support \\
\cmd{muse-py serve}, \cmd{dicom-py serve} & \code{VIEWER\_ROOT=DIR} when the package is installed
  (default: the whole file system) & The most that \cmd{serve} ever serves: a \opt{-{}-root} that
  contains it is narrowed to it, and one outside it is refused \\
\cmd{muse-py serve}, \cmd{dicom-py serve} & \code{VIEWER\_PASSWORD=0} (default \code{1}) & No
  password support \\
The C++ project & CMake \code{VS\_BUILD\_VIEWER} (default \code{ON}) & The viewer. Off: only the
  libraries \code{vs\_core} and \code{vs\_engine}, for adapters written in C++ \\
The C++ project & CMake \code{VS\_BUILD\_TESTS} (default \code{ON}) & The unit tests
  (GoogleTest) \\
The C++ project & CMake \code{VS\_DEPS\_DIR} (default \file{cpp/deps}) & The folder of
  pre-downloaded dependency archives, for builds without internet access \\
\end{xltabular}

At run time, the settings (\code{data\_roots}, \opt{-{}-root}) can narrow a viewer's root
further, but never widen it. The standalone installer builds in folders of its own
(\file{target/standalone}, \file{viewers/vitalstream/cpp/build/standalone}), so that Medorama's
viewer in the same clone (\file{target/release}, \file{build/viewer}) keeps its own options: the
whole file system as its root, and password support.

\begin{xltabular}{\linewidth}{@{}P{2.4cm}P{3.4cm}L@{}}
\caption{The C++ project's CMake presets.}
\label{tab:settings-reference-presets}\\
\toprule
\tabhead Preset & \tabhead Folder & \tabhead Builds \\
\midrule
\endfirsthead
\toprule
\tabhead Preset & \tabhead Folder & \tabhead Builds \\
\midrule
\endhead
\bottomrule
\endfoot
\code{release} & \file{build/release} & The viewer and its tests \\
\code{debug} & \file{build/debug} & The same, unoptimized, with debug information \\
\code{viewer} & \file{build/viewer} & The viewer only (what the installers use) \\
\end{xltabular}
